%% ============================================================================
%%  LAFROX_AnalisisFinanciero_2_Anexos_20261113.tex -- Anexo A del Análisis económico-financiero: tabla de trazabilidad
%%  observación–respuesta–sección modificada de las observaciones del Informe 2
%%  (Art. 46 de las Bases Administrativas).
%%
%%  Las tablas se generan en partes/tablas.tex y partes/tablas_sintesis.tex.
%%  Compilar:  latexmk LAFROX_AnalisisFinanciero_2_Anexos_20261113.tex   (LuaLaTeX, ver latexmkrc)
%% ============================================================================
\documentclass[carta,firma]{lafrox}
\usepackage{needspace}
\usepackage{xurl}

\datoslafrox{
  caso         = {Caso --- Logística},
  documento    = {Oferta Económica},
  sobre        = {Sobre N.\textdegree{} 3 \textemdash{} Oferta Económica},
  fecha        = {13 de noviembre de 2026},
  subtitulo    = {Análisis económico-financiero --- Anexos},
  alcance      = {Entregable 2},
  formulario   = {Anexo A --- Tabla de trazabilidad (Art. 46)},
  nomenclatura = {LAFROX\_AnalisisFinanciero\_2\_Anexos\_20261113.pdf},
}

\configurarportadalafrox{Análisis económico-financiero \textemdash{} Anexos}{Anexo A \textemdash{} Trazabilidad de observaciones del Informe 2}{Art. 46}{}{FIRMA_Alex_Aravena_transparente.png}
\makeatletter\def\LFXfecha{13 de noviembre de 2026}\makeatother

\AtBeginDocument{\hypersetup{pdftitle={Anexo A: Trazabilidad de observaciones del Informe 2},pdfauthor={LafroX SpA}}}
\renewcommand{\thesection}{A.\arabic{section}}
\renewcommand{\thesubsection}{\thesection.\arabic{subsection}}
\renewcommand{\thetable}{\arabic{table}}
\setcounter{tocdepth}{2}

\begin{document}
\portadalafrox
\indicelafrox

\Needspace{6\baselineskip}
\section{Alcance y criterio de lectura}

Este anexo responde el Art. 46 de las Bases Administrativas, que exige resolver en forma explícita las observaciones de la instancia anterior mediante una tabla observación–respuesta–sección modificada (Distribuidora Puelche S.A., 2026a, Art. 46, p. 28). Reúne las observaciones de la primera revisión del Informe 2 a los Subdocumentos 1 a 9 y 13 y a sus Formularios T-6, T-9 a T-19, más las observaciones generales de formalidad.

Se registran las observaciones que piden un cambio: las marcadas como «Parcial», «No corregida» o «Crítico» y las que la revisión formula sin calificativo. Los aspectos que la revisión calificó como «OK» o «Corregida» no se repiten. Cuando una misma observación reaparece en «Qué se espera en la Propuesta Final» o en las consideraciones transversales, se registra una sola vez, en el subdocumento donde se corrige.

Cada respuesta comienza con su estado. «Corregida» indica que el cambio está en la versión vigente del subdocumento o formulario. «Parcial», que el cambio atiende una parte de la observación y lo que falta se completa para la Propuesta Final. «Pendiente», que la corrección se incorpora en la Propuesta Final. La columna «Sección modificada» indica el lugar donde se lee la respuesta. El identificador combina el subdocumento y un número correlativo: SD3-04 es la cuarta observación registrada del Subdocumento 3; SD4.1 y SD4.2 corresponden a las secciones 4.1 y 4.2–4.3 del Subdocumento 4, y GEN a las observaciones generales.

La Tabla~\ref{tab:sintesis} resume el estado por subdocumento y las Tablas 2 a 13 presentan el detalle. La mayor parte de las observaciones quedó corregida en los subdocumentos que venían del Informe 1 y en los Subdocumentos 6, 8 y 9. Las parciales y pendientes, concentradas en el Subdocumento 7 y en puntos acotados de la formalidad general y de los Subdocumentos 3 y 4, quedan resueltas para la Propuesta Final.

\Needspace{20\baselineskip}
\begin{tablalafrox}{Síntesis del estado de las observaciones por subdocumento. Fuente: Tablas 2 a 13 de este anexo.}{tab:sintesis}{L{6.2cm} C{1.9cm} C{1.6cm} C{1.9cm} C{1.4cm}}
  {\cab{Subdocumento} & \cab{Corregidas} & \cab{Parciales} & \cab{Pendientes} & \cab{Total}}
  General & 3 & 1 & 2 & 6 \\
  Subdocumento 1 & 9 & 0 & 0 & 9 \\
  Subdocumento 2 & 7 & 0 & 0 & 7 \\
  Subdocumento 3 & 7 & 5 & 2 & 14 \\
  Subdocumento 4.1 & 9 & 0 & 0 & 9 \\
  Subdocumento 4.2 y 4.3 & 7 & 1 & 1 & 9 \\
  Subdocumento 5 & 7 & 2 & 0 & 9 \\
  Subdocumento 6 & 16 & 1 & 0 & 17 \\
  Subdocumento 7 & 12 & 9 & 5 & 26 \\
  Subdocumento 8 & 12 & 7 & 2 & 21 \\
  Subdocumento 9 & 8 & 1 & 0 & 9 \\
  Subdocumento 13 & 3 & 3 & 0 & 6 \\
  \textbf{Total} & \textbf{100} & \textbf{30} & \textbf{12} & \textbf{142} \\
\end{tablalafrox}

\clearpage
\Needspace{6\baselineskip}
\section{Observaciones, respuestas y secciones modificadas}
\Needspace{8\baselineskip}
\subsection{General: formalidad y cumplimiento de instrucciones}
{\linespread{1}\LFXcuerpotabla
\setlength{\parskip}{0pt}\renewcommand{\arraystretch}{1.10}\setlength{\tabcolsep}{4pt}
\begin{longtable}{F{0.09} F{0.29} F{0.40} F{0.22}}
\caption[General: formalidad y cumplimiento de instrucciones]{General: formalidad y cumplimiento de instrucciones. Fuente: revisión del Informe 2 y versión corregida de LafroX.}\\[2pt]
\LFXtopruleoriginal \cab{N.º} & \cab{Observación} & \cab{Respuesta} & \cab{Sección modificada} \\ \midrule
\endfirsthead
\multicolumn{4}{@{}l@{}}{\fontsize{9pt}{11pt}\selectfont\color{lafroxGris}Tabla \thetable\ (continuación)}\\[3pt]
\LFXtopruleoriginal \cab{N.º} & \cab{Observación} & \cab{Respuesta} & \cab{Sección modificada} \\ \midrule
\endhead
\midrule
\multicolumn{4}{@{}r@{}}{\fontsize{9pt}{11pt}\selectfont\color{lafroxGris}Continúa en el folio siguiente.}\\
\endfoot
\bottomrule
\endlastfoot
\mbox{GEN-01} & Volumen excesivo: SD4 con 194 páginas de cuerpo y 117 de anexos; el subdocumento debe resumir y el anexo detallar. & \textbf{Parcial.} El SD4 traslada al Anexo 4-D los listados de M2, M3, M4 y M7, al 4-N los criterios de interacción y al 4-R la custodia. La reducción del cuerpo del SD4 y del SD7 continúa en la Propuesta Final. & SD4, 4.1.5 y Anexos 4-D, 4-N, 4-R; SD7 \\
\LFXfiletefila
\mbox{GEN-02} & La tabla del Art. 46° llega en un archivo fuera de la nomenclatura de la sección 2 de la Aclaración N.° 2. & \textbf{Corregida.} La tabla se entrega como Anexo A del Análisis económico-financiero (Entregable 2), en archivo propio con nombre conforme a la nomenclatura de la oferta. & Este anexo \\
\LFXfiletefila
\mbox{GEN-03} & Declaración de uso de IA con nivel «Alto» en todo y un único revisor; revisión humana de forma. & \textbf{Corregida.} Las Declaraciones de SD1, SD3, SD4, SD5, SD6 y SD8 asignan revisores distintos por sección, cada uno con el contenido que verificó; el SD13 agrega una pauta de revisión técnica. & Declaraciones de IA de SD1, SD3, SD4, SD5, SD6 y SD8; Anexos SD13 \\
\LFXfiletefila
\mbox{GEN-04} & Foliación única y correlativa de toda la propuesta. & \textbf{Pendiente.} La foliación correlativa, de la portada del SD1 al último formulario, se aplica al ensamblar la Propuesta Final. & Ensamble de la Propuesta Final \\
\LFXfiletefila
\mbox{GEN-05} & Erratas dentro de las figuras del SD2 («enERP», «enpapel», «retornarr», «No Diferencias no conciliadas»). & \textbf{Corregida.} Figuras reexportadas desde su fuente editable con el texto corregido; «No» queda como rótulo del conector. & SD2, Figuras 2.1, 2.5 y 2.6 \\
\LFXfiletefila
\mbox{GEN-06} & Sitio web corporativo en GitHub (lafrox.github.io) y correo con dominio lafrox.cl. & \textbf{Pendiente.} El dominio del sitio corporativo se unifica con el del correo. Se corrige para la Propuesta Final. & SD1, introducción \\
\end{longtable}
}

\Needspace{8\baselineskip}
\subsection{Subdocumento 1: Presentación de la empresa (T-6)}
{\linespread{1}\LFXcuerpotabla
\setlength{\parskip}{0pt}\renewcommand{\arraystretch}{1.10}\setlength{\tabcolsep}{4pt}
\begin{longtable}{F{0.09} F{0.29} F{0.40} F{0.22}}
\caption[Subdocumento 1: Presentación de la empresa (T-6)]{Subdocumento 1: Presentación de la empresa (T-6). Fuente: revisión del Informe 2 y versión corregida de LafroX.}\\[2pt]
\LFXtopruleoriginal \cab{N.º} & \cab{Observación} & \cab{Respuesta} & \cab{Sección modificada} \\ \midrule
\endfirsthead
\multicolumn{4}{@{}l@{}}{\fontsize{9pt}{11pt}\selectfont\color{lafroxGris}Tabla \thetable\ (continuación)}\\[3pt]
\LFXtopruleoriginal \cab{N.º} & \cab{Observación} & \cab{Respuesta} & \cab{Sección modificada} \\ \midrule
\endhead
\midrule
\multicolumn{4}{@{}r@{}}{\fontsize{9pt}{11pt}\selectfont\color{lafroxGris}Continúa en el folio siguiente.}\\
\endfoot
\bottomrule
\endlastfoot
\mbox{SD1-01} & Sin antecedentes financieros: ventas, patrimonio ni capacidad de garantía. & \textbf{Corregida.} Se agregan patrimonio 2023–2025, ventas 2025, liquidez corriente de 1,8 y capacidad de emisión de garantías. La acreditación documental se entrega en el Sobre 1. & SD1, 1.1 \\
\LFXfiletefila
\mbox{SD1-02} & El SOC aparece como existente, contratado y por implementar. & \textbf{Corregida.} Una sola versión: NOC y función CISO existentes; SOC propio en Valparaíso por habilitar con personal adicional, con calendario en el T-15. & SD1, 1.1, 1.2 y 1.5 \\
\LFXfiletefila
\mbox{SD1-03} & Demanda de desarrollo poco clara (31 FTE frente a 40 simultáneos y 48 disponibles); debe cuadrar con el T-15. & \textbf{Corregida.} Se retiran las cifras incompatibles y queda una capacidad única: 48 desarrolladores y 15 especialistas de arquitectura e integración. El T-15 muestra un máximo de 32 plazas de Desarrollo cubiertas por 33 personas. & SD1, 1.2 y 1.5; T-15, 5.7 \\
\LFXfiletefila
\mbox{SD1-04} & El organigrama omite Comercial y Administración y no ubica los dos centros de datos. & \textbf{Corregida.} Nueva Figura 1.1 con Comercial y Administración y Finanzas bajo la Gerencia; centro primario en Santiago y secundario en Valparaíso; se agrega el área de Implantación y Gestión del Cambio (9 especialistas, dentro de una dotación técnica de 133). & SD1, 1.2, Figura 1.1 y Tabla 1.3 \\
\LFXfiletefila
\mbox{SD1-05} & La sección 1.5 nomina al equipo del proyecto, reservado al Capítulo 12. & \textbf{Corregida.} La 1.5 describe siete ámbitos y diez roles sin nombres y remite el equipo nominado al T-8. & SD1, 1.5 \\
\LFXfiletefila
\mbox{SD1-06} & El plan de renovación de certificaciones se presenta dos veces. & \textbf{Corregida.} Se elimina el párrafo repetido y se conservan las fichas por certificación; se corrige el plazo atribuido a ISO/IEC 20000-1. & SD1, 1.4 \\
\LFXfiletefila
\mbox{SD1-07} & No acredita el conocimiento sanitario y tributario de la RT-15.02 del caso. & \textbf{Corregida.} Se declara el conocimiento del Reglamento Sanitario de los Alimentos y de los DTE, respaldado en registros de capacitación y en los proyectos del T-6. & SD1, 1.4 \\
\LFXfiletefila
\mbox{SD1-08} & Las cuatro alianzas vencen entre los meses 2 y 11; falta analizarlo. & \textbf{Corregida.} Análisis del primer año: renovación 60 días antes del vencimiento, separación entre membresía y soporte y contingencia. & SD1, 1.6 \\
\LFXfiletefila
\mbox{SD1-09} & Referencia farmacéutica con cámaras a −22 °C poco habitual. & \textbf{Corregida.} Se precisa que son doce cámaras de reactivos y materiales de diagnóstico congelados. & SD1, 1.4, Tabla 1.5; T-6, proyecto 2 \\
\end{longtable}
}

\Needspace{8\baselineskip}
\subsection{Subdocumento 2: Comprensión del problema}
{\linespread{1}\LFXcuerpotabla
\setlength{\parskip}{0pt}\renewcommand{\arraystretch}{1.10}\setlength{\tabcolsep}{4pt}
\begin{longtable}{F{0.09} F{0.29} F{0.40} F{0.22}}
\caption[Subdocumento 2: Comprensión del problema]{Subdocumento 2: Comprensión del problema. Fuente: revisión del Informe 2 y versión corregida de LafroX.}\\[2pt]
\LFXtopruleoriginal \cab{N.º} & \cab{Observación} & \cab{Respuesta} & \cab{Sección modificada} \\ \midrule
\endfirsthead
\multicolumn{4}{@{}l@{}}{\fontsize{9pt}{11pt}\selectfont\color{lafroxGris}Tabla \thetable\ (continuación)}\\[3pt]
\LFXtopruleoriginal \cab{N.º} & \cab{Observación} & \cab{Respuesta} & \cab{Sección modificada} \\ \midrule
\endhead
\midrule
\multicolumn{4}{@{}r@{}}{\fontsize{9pt}{11pt}\selectfont\color{lafroxGris}Continúa en el folio siguiente.}\\
\endfoot
\bottomrule
\endlastfoot
\mbox{SD2-01} & El arbitraje ya adopta decisiones de solución (promesa 24/48 h con corte a las 14:00, retención por SSCC). & \textbf{Corregida.} El SD2 deja la tensión y su condición; el corte, los plazos y el tratamiento térmico pasan al SD3 (3.2.2, S-21, D-03 y RNG-15). & SD2, 2.2.3 y 2.5.1; SD3, 3.2.2 \\
\LFXfiletefila
\mbox{SD2-02} & La Figura 2.3 omite la preparación nocturna con hoja de picking y el faltante sin causa. & \textbf{Corregida.} La figura pasa a dos paneles: circuito de la hoja de picking y pérdida de la causa del faltante, distinguiendo los despachos directos. & SD2, 2.3 (proceso 3), Figura 2.3 \\
\LFXfiletefila
\mbox{SD2-03} & La Tabla 2.3 sintetiza 5 de las 12 restricciones y omite las operacionales. & \textbf{Corregida.} Se reemplaza la tabla por cuatro párrafos que cubren R-01 a R-12. & SD2, 2.5.3 \\
\LFXfiletefila
\mbox{SD2-04} & Afirmaciones que el caso no hace (WMS con soporte discontinuado; módulos del ERP abandonados). & \textbf{Corregida.} Se distingue la aplicación de preventa de 2016 sin proveedor del ERP de 2017 y se describe la cobertura real del WMS. & SD2, 2.1 \\
\LFXfiletefila
\mbox{SD2-05} & La rotación del 38 \% aparece dentro del costo de servir. & \textbf{Corregida.} Se traslada a la preparación nocturna, sin atribuirle una tasa de error no medida. & SD2, 2.2.1 y 2.3 \\
\LFXfiletefila
\mbox{SD2-06} & Cita de Boulaksil y Belkora: Casablanca en el texto y Valencia en la referencia. & \textbf{Corregida.} Se aclara que Valencia es el fabricante estudiado en Casablanca y se acota la atribución al artículo. & SD2, 2.2.4; Referencias \\
\LFXfiletefila
\mbox{SD2-07} & La Declaración de IA atribuye a la fila 2.1 la validación de la propuesta de valor. & \textbf{Corregida.} Se reemplaza por el registro de revisión del resumen ejecutivo. & SD2, Tabla 2.3 (Declaración de IA) \\
\end{longtable}
}

\Needspace{8\baselineskip}
\subsection{Subdocumento 3: Esquema de solución y alcance (T-12)}
{\linespread{1}\LFXcuerpotabla
\setlength{\parskip}{0pt}\renewcommand{\arraystretch}{1.10}\setlength{\tabcolsep}{4pt}
\begin{longtable}{F{0.09} F{0.29} F{0.40} F{0.22}}
\caption[Subdocumento 3: Esquema de solución y alcance (T-12)]{Subdocumento 3: Esquema de solución y alcance (T-12). Fuente: revisión del Informe 2 y versión corregida de LafroX.}\\[2pt]
\LFXtopruleoriginal \cab{N.º} & \cab{Observación} & \cab{Respuesta} & \cab{Sección modificada} \\ \midrule
\endfirsthead
\multicolumn{4}{@{}l@{}}{\fontsize{9pt}{11pt}\selectfont\color{lafroxGris}Tabla \thetable\ (continuación)}\\[3pt]
\LFXtopruleoriginal \cab{N.º} & \cab{Observación} & \cab{Respuesta} & \cab{Sección modificada} \\ \midrule
\endhead
\midrule
\multicolumn{4}{@{}r@{}}{\fontsize{9pt}{11pt}\selectfont\color{lafroxGris}Continúa en el folio siguiente.}\\
\endfoot
\bottomrule
\endlastfoot
\mbox{SD3-01} & Figura 3.4 ilegible; figuras delgadas (la 3.2 son cuatro cajas) y ninguna muestra qué corre en el CD, plataformas, camión y nube. & \textbf{Parcial.} Se agrega en 3.4.2 la síntesis de emplazamiento por sitio, camión y nube (Anexo 4-N). La Figura 3.4 aún no se redibuja. Lo que falta se completa para la Propuesta Final. & SD3, 3.4.2 \\
\LFXfiletefila
\mbox{SD3-02} & Erratas en figuras: «Gerente de operacions» y rótulos invertidos en la Figura 3.7. & \textbf{Pendiente.} Las Figuras 3.4, 3.5 y 3.7 se reeditan con los rótulos corregidos para la Propuesta Final. & SD3, Figuras 3.4, 3.5 y 3.7 \\
\LFXfiletefila
\mbox{SD3-03} & Columna «Origen» de la Tabla 3.A.1 vacía o cortada. & \textbf{Corregida.} Se reconstruye el origen de todos los RF y se precisan RNF y requisitos propios; se retiran los orígenes genéricos. & SD3-Anexos, Anexo 3.A \\
\LFXfiletefila
\mbox{SD3-04} & La prueba y el criterio de aceptación quedan fuera del T-12 (cadena del 17.1). & \textbf{Corregida.} El T-12 se enlaza con la matriz única del T-13 (origen, requisito, componente, EDT, prueba y criterio); el T-17 conserva el criterio contractual. & T-12; T-13, matriz de requisitos \\
\LFXfiletefila
\mbox{SD3-05} & Prioridades poco discriminantes (84 \% en Crítica o Alta). & \textbf{Parcial.} Se aplica el techo por familia: 58 críticos, 157 altos y 55 medios, y se explica la diferencia entre prioridad y etapa. Lo que falta se completa para la Propuesta Final. & SD3, 3.2; Anexo 3.A \\
\LFXfiletefila
\mbox{SD3-06} & La columna «Componente» del T-12 repite la pila completa en cada fila. & \textbf{Pendiente.} Se precisaron componentes en filas puntuales, sin retirar la repetición. Se corrige para la Propuesta Final. & T-12 \\
\LFXfiletefila
\mbox{SD3-07} & Interoperabilidad, escalabilidad, desempeño y seguridad en cuatro viñetas. & \textbf{Parcial.} La síntesis de emplazamiento complementa el análisis; el detalle sigue en el SD4. Lo que falta se completa para la Propuesta Final. & SD3, 3.4.2 \\
\LFXfiletefila
\mbox{SD3-08} & Anexo 3.F reducido; fundamentos que obligan a leer SD2 y SD4. & \textbf{Parcial.} Se amplía el fundamento de D-03 junto con S-21 y RNG-15. Lo que falta se completa para la Propuesta Final. & SD3-Anexos, Anexo 3.F \\
\LFXfiletefila
\mbox{SD3-09} & Nombres de módulos distintos en SD3 y SD4 (M10 «BI» y M11 «EDI»). & \textbf{Corregida.} Las figuras del SD4 usan los nombres del SD3: M10 Analítica y M11 Canal moderno. & SD4, Figura 1 y variantes; Anexo 4-D \\
\LFXfiletefila
\mbox{SD3-10} & Meta de OTIF del mes 32 contradictoria entre el cuerpo y el Anexo 3.J. & \textbf{Corregida.} Una sola regla: estrictamente mayor que 95 \%; 90 \% en el mes 15 y 93 \% en el mes 19. & SD3, 3.3; Anexo 3.J; T-17, R18-04 \\
\LFXfiletefila
\mbox{SD3-11} & Incorporación de conductores externos despachada en una frase. & \textbf{Corregida.} Se describe la declaración del transportista, el alta con código de un solo uso, la aprobación de Despacho y la reasignación a un suplente. & SD3, 3.4.4 \\
\LFXfiletefila
\mbox{SD3-12} & La operación no cuantifica la dotación de la mesa ni del NOC. & \textbf{Corregida.} Se remiten expresamente al Subdocumento 10 los turnos, guardias y refuerzos. & SD3, 3.4.5 \\
\LFXfiletefila
\mbox{SD3-13} & Revisión humana declarada sólo de forma; 3.4 declara «Ninguno» en diagramas. & \textbf{Parcial.} Revisores distintos por cuerpo y anexos, con el contenido cotejado. La fila 3.4 mantiene «Ninguno» en diagramas. Lo que falta se completa para la Propuesta Final. & SD3, Tabla 3.6 \\
\LFXfiletefila
\mbox{SD3-14} & El SD3 no menciona ninguna de las cinco innovaciones. & \textbf{Corregida.} Se incorporan las cinco innovaciones al alcance, con remisión al SD13 y al T-19. & SD3, 3.1 y 3.4 \\
\end{longtable}
}

\Needspace{8\baselineskip}
\subsection{Subdocumento 4.1: Arquitectura lógica}
{\linespread{1}\LFXcuerpotabla
\setlength{\parskip}{0pt}\renewcommand{\arraystretch}{1.10}\setlength{\tabcolsep}{4pt}
\begin{longtable}{F{0.09} F{0.29} F{0.40} F{0.22}}
\caption[Subdocumento 4.1: Arquitectura lógica]{Subdocumento 4.1: Arquitectura lógica. Fuente: revisión del Informe 2 y versión corregida de LafroX.}\\[2pt]
\LFXtopruleoriginal \cab{N.º} & \cab{Observación} & \cab{Respuesta} & \cab{Sección modificada} \\ \midrule
\endfirsthead
\multicolumn{4}{@{}l@{}}{\fontsize{9pt}{11pt}\selectfont\color{lafroxGris}Tabla \thetable\ (continuación)}\\[3pt]
\LFXtopruleoriginal \cab{N.º} & \cab{Observación} & \cab{Respuesta} & \cab{Sección modificada} \\ \midrule
\endhead
\midrule
\multicolumn{4}{@{}r@{}}{\fontsize{9pt}{11pt}\selectfont\color{lafroxGris}Continúa en el folio siguiente.}\\
\endfoot
\bottomrule
\endlastfoot
\mbox{SD4.1-01} & Figura 1 (vista general) ilegible a tamaño impreso. & \textbf{Corregida.} Nueva vista vectorial de navegación por ocho capas y doce módulos, con letra mínima de 10,4 pt. & SD4, 4.1, Figura 1 \\
\LFXfiletefila
\mbox{SD4.1-02} & Matriz requerimiento → módulo → capa por grupo y no por requisito. & \textbf{Corregida.} Dieciséis obligaciones individuales con origen, componente, capa, contrato, prueba y aceptación, y cuatro recorridos críticos en el cuerpo. & SD4, Tabla 4.1-A; Anexo 4-E, Tabla A.6a \\
\LFXfiletefila
\mbox{SD4.1-03} & Leyendas de fuente de plantilla, iguales en todas las figuras. & \textbf{Corregida.} Cada figura lógica y física cita el dato que la sostiene. & SD4, leyendas de 4.1, 4.2 y 4.3; Anexos \\
\LFXfiletefila
\mbox{SD4.1-04} & Cambio de Django a Laravel sin reconocerlo frente al Informe 1. & \textbf{Corregida.} Se explica el antecedente Django, el motivo de la elección y sus efectos en contratos, pruebas y herramientas. & SD4, 4.1.7.1; ADR-23 \\
\LFXfiletefila
\mbox{SD4.1-05} & Versiones recientes (Laravel 13, PHP 8.5, Angular 22, TypeScript 6.0) por demostrar. & \textbf{Corregida.} Tabla de lanzamiento, restricciones y soporte de ocho productos con fuentes oficiales y política de actualización. & SD4, Anexo 4-P, Tabla A.20 \\
\LFXfiletefila
\mbox{SD4.1-06} & Referencia a la «sección 2.5» inexistente del caso. & \textbf{Corregida.} Se cita la clasificación correcta (RT-06.01). & SD4, 4.3.1 \\
\LFXfiletefila
\mbox{SD4.1-07} & Párrafo con redacción defectuosa sobre disponibilidad y redundancia. & \textbf{Corregida.} Se reescribe el párrafo. & SD4, 4.3.1 \\
\LFXfiletefila
\mbox{SD4.1-08} & Credencial de turno de 14 h frente a rutas de hasta 14,5 h; manifiesto de 26 h. & \textbf{Corregida.} Asignación emitida dentro de la hora previa a salir, válida hasta el fin programado más 30 min y con máximo de 16 h; manifiesto con vigencia remanente mínima de 25 h y dos relevos. & SD4, 4.1 (identidad y conductores externos); ADR-06; Anexo 4-M \\
\LFXfiletefila
\mbox{SD4.1-09} & Declaración de IA «Alto» en texto y diagramas con un solo revisor; 4.3 declara «Ninguno». & \textbf{Corregida.} Cinco revisores se reparten las secciones, anexos y el T-11, cada uno con el contenido que verificó. & SD4, Tabla 29 (Declaración de IA) \\
\end{longtable}
}

\Needspace{8\baselineskip}
\subsection{Subdocumento 4.2 y 4.3: Arquitectura física y centros de datos (T-11)}
{\linespread{1}\LFXcuerpotabla
\setlength{\parskip}{0pt}\renewcommand{\arraystretch}{1.10}\setlength{\tabcolsep}{4pt}
\begin{longtable}{F{0.09} F{0.29} F{0.40} F{0.22}}
\caption[Subdocumento 4.2 y 4.3: Arquitectura física y centros de datos (T-11)]{Subdocumento 4.2 y 4.3: Arquitectura física y centros de datos (T-11). Fuente: revisión del Informe 2 y versión corregida de LafroX.}\\[2pt]
\LFXtopruleoriginal \cab{N.º} & \cab{Observación} & \cab{Respuesta} & \cab{Sección modificada} \\ \midrule
\endfirsthead
\multicolumn{4}{@{}l@{}}{\fontsize{9pt}{11pt}\selectfont\color{lafroxGris}Tabla \thetable\ (continuación)}\\[3pt]
\LFXtopruleoriginal \cab{N.º} & \cab{Observación} & \cab{Respuesta} & \cab{Sección modificada} \\ \midrule
\endhead
\midrule
\multicolumn{4}{@{}r@{}}{\fontsize{9pt}{11pt}\selectfont\color{lafroxGris}Continúa en el folio siguiente.}\\
\endfoot
\bottomrule
\endlastfoot
\mbox{SD4.2-01} & La sala no declara superficie ni tiene plano acotado (RT-06.03); recintos de más para dos racks. & \textbf{Corregida.} Se declara la superficie útil del recinto y la Figura 29 pasa a plano acotado por zonas. & SD4, 4.3.1, tabla de superficie y Figura 29 \\
\LFXfiletefila
\mbox{SD4.2-02} & Sin alternativa si el CLIENTE no aprueba la residencia en sa-east-1 y us-east-1. & \textbf{Pendiente.} El ADR-19 mantiene la aprobación del CLIENTE como hito de la Etapa 1, sin arquitectura alternativa. Se corrige para la Propuesta Final. & SD4, Anexo 4-O, ADR-19 \\
\LFXfiletefila
\mbox{SD4.2-03} & Estudio de cobertura dentro de las cámaras diferido. & \textbf{Corregida.} Estudio con diseño predictivo por zona y verificación en sitio en el mes 2 (paquete 2.1.2). & SD4, Anexo 4-W.13; T-14, 2.1.2 \\
\LFXfiletefila
\mbox{SD4.2-04} & Figura 18 (topología AWS) con rótulos de 6 a 7 pt. & \textbf{Corregida.} Figura redibujada a 9 pt. & SD4, 4.2, Figura 18 \\
\LFXfiletefila
\mbox{SD4.2-05} & La Tabla 14 no explica el total por máximo coincidente. & \textbf{Corregida.} Nota que explica que el total es el máximo agregado coincidente y no la suma de máximos. & SD4, 4.2.6, Tabla 14 \\
\LFXfiletefila
\mbox{SD4.2-06} & LafroX provee la cadena de frío, contra la exclusión del Cap. 11 del caso. & \textbf{Corregida.} El CLIENTE adquiere los sensores y kits según la especificación del T-11; LafroX los instala, integra y mantiene. Alineado en SD3, SD6, T-9 y SD7. & SD4, 4.2.3; T-11; SD6, 6.1; SD7, 5.1.4 \\
\LFXfiletefila
\mbox{SD4.2-07} & El T-11 no declara ciclo de vida ni plan de reposición por equipo. & \textbf{Corregida.} Tabla de ciclo de vida, repuestos y reposición por dispositivo en los 56 meses. & T-11, ciclo de vida por dispositivo \\
\LFXfiletefila
\mbox{SD4.2-08} & Sobra la frase «los niveles del 7.2 son un medio, no un fin». & \textbf{Corregida.} Frase eliminada. & SD4, 4.3 \\
\LFXfiletefila
\mbox{SD4.2-09} & Riesgo regulatorio y de latencia de la residencia fuera de Chile no llevado al plan de riesgos. & \textbf{Parcial.} El T-16 registra la falta de aprobación de regiones de nube como riesgo vigente; la latencia no se cuantifica. Lo que falta se completa para la Propuesta Final. & T-16; SD8, Anexo 8.F \\
\end{longtable}
}

\Needspace{8\baselineskip}
\subsection{Subdocumento 5: Modelo y gestión de datos}
{\linespread{1}\LFXcuerpotabla
\setlength{\parskip}{0pt}\renewcommand{\arraystretch}{1.10}\setlength{\tabcolsep}{4pt}
\begin{longtable}{F{0.09} F{0.29} F{0.40} F{0.22}}
\caption[Subdocumento 5: Modelo y gestión de datos]{Subdocumento 5: Modelo y gestión de datos. Fuente: revisión del Informe 2 y versión corregida de LafroX.}\\[2pt]
\LFXtopruleoriginal \cab{N.º} & \cab{Observación} & \cab{Respuesta} & \cab{Sección modificada} \\ \midrule
\endfirsthead
\multicolumn{4}{@{}l@{}}{\fontsize{9pt}{11pt}\selectfont\color{lafroxGris}Tabla \thetable\ (continuación)}\\[3pt]
\LFXtopruleoriginal \cab{N.º} & \cab{Observación} & \cab{Respuesta} & \cab{Sección modificada} \\ \midrule
\endhead
\midrule
\multicolumn{4}{@{}r@{}}{\fontsize{9pt}{11pt}\selectfont\color{lafroxGris}Continúa en el folio siguiente.}\\
\endfoot
\bottomrule
\endlastfoot
\mbox{SD5-01} & Figura 5.1 ilegible (texto de 1 a 2 pt). & \textbf{Corregida.} Se reemplaza por una vista de dominios y relaciones principales; el detalle queda en las Figuras 5.2 a 5.10. & SD5, 5.1, Figura 5.1 \\
\LFXfiletefila
\mbox{SD5-02} & Vistas por dominio cerca de 8 pt y título repetido en la imagen. & \textbf{Parcial.} Figuras reexportadas y ampliadas a la altura útil de la página; falta comprobar los 9 pt en todas. Lo que falta se completa para la Propuesta Final. & SD5, Figuras 5.2 a 5.10 \\
\LFXfiletefila
\mbox{SD5-03} & Contradicción sobre «M» entre el cuerpo y el glosario. & \textbf{Corregida.} «M» queda reservada a módulos; los plazos se escriben «mes N». & SD5, introducción; Anexo 5-A, glosario \\
\LFXfiletefila
\mbox{SD5-04} & No se declara de dónde sale el umbral de 0,3 \% de desviación de inventario. & \textbf{Corregida.} Se fundamenta como control de riesgo de LafroX para la carga inicial. & SD5, 5.3 \\
\LFXfiletefila
\mbox{SD5-05} & El SD5 no muestra el cálculo de los volúmenes a migrar. & \textbf{Corregida.} Cálculo por conjunto (registros × meses × KB): 16.054.780 KB de origen y 32,1 GB de destino, con sensibilidad. & SD5, 5.3 \\
\LFXfiletefila
\mbox{SD5-06} & Prosa telegráfica, cargada de barras, siglas y remisiones. & \textbf{Parcial.} Se reescriben en prosa explicativa la introducción, 5.1, 5.3 y 5.4; parte de 5.2 sigue comprimida. Lo que falta se completa para la Propuesta Final. & SD5, 5.1, 5.3 y 5.4 \\
\LFXfiletefila
\mbox{SD5-07} & Dos nomenclaturas para los subdocumentos (SD1–SD5 y SD1–SD5) y nombre del archivo en la portada. & \textbf{Corregida.} Se usa sólo «SD» y se retira el nombre del archivo de la portada. & SD5, cuerpo y portada \\
\LFXfiletefila
\mbox{SD5-08} & Declaración de IA «Alto» con un solo revisor. & \textbf{Corregida.} Cinco revisores se reparten las secciones y anexos, cada uno con el contenido que verificó. & SD5, Declaración de IA \\
\LFXfiletefila
\mbox{SD5-09} & Los relevos de 8/16 h y la credencial de turno no se explican en el SD5. & \textbf{Corregida.} Se explica la función de cada plazo: identidad de consulta de 24 h, credencial del turno y manifiesto de 26 h para relevos inscritos, con su prueba de corte y cambio de turno. & SD5, 5.2, Tabla 5.1 \\
\end{longtable}
}

\Needspace{8\baselineskip}
\subsection{Subdocumento 6: Metodologías (T-9 y T-10)}
{\linespread{1}\LFXcuerpotabla
\setlength{\parskip}{0pt}\renewcommand{\arraystretch}{1.10}\setlength{\tabcolsep}{4pt}
\begin{longtable}{F{0.09} F{0.29} F{0.40} F{0.22}}
\caption[Subdocumento 6: Metodologías (T-9 y T-10)]{Subdocumento 6: Metodologías (T-9 y T-10). Fuente: revisión del Informe 2 y versión corregida de LafroX.}\\[2pt]
\LFXtopruleoriginal \cab{N.º} & \cab{Observación} & \cab{Respuesta} & \cab{Sección modificada} \\ \midrule
\endfirsthead
\multicolumn{4}{@{}l@{}}{\fontsize{9pt}{11pt}\selectfont\color{lafroxGris}Tabla \thetable\ (continuación)}\\[3pt]
\LFXtopruleoriginal \cab{N.º} & \cab{Observación} & \cab{Respuesta} & \cab{Sección modificada} \\ \midrule
\endhead
\midrule
\multicolumn{4}{@{}r@{}}{\fontsize{9pt}{11pt}\selectfont\color{lafroxGris}Continúa en el folio siguiente.}\\
\endfoot
\bottomrule
\endlastfoot
\mbox{SD6-01} & Regla de valor ganado truncada («una desviación mayor al 10»). & \textbf{Corregida.} Regla completa en el cuerpo y en el T-9, con tabla de umbrales: SPI o CPI bajo 0,90 o desviación de más de 10 \% en la ruta crítica obliga a un plan de recuperación en cinco días hábiles, aprobado por el Comité de Proyecto. & SD6, 6.1.5; T-9, 1.5, Tabla T-9.4 \\
\LFXfiletefila
\mbox{SD6-02} & PMBOK «adaptado» sin explicar la adaptación ni la convivencia de RUP, Kanban e iteraciones con los hitos. & \textbf{Corregida.} Frontera entre lo que se fija y lo que se adapta, áreas del PMBOK reforzadas y aligeradas, y modelo en tres niveles (fases RUP e hitos, iteraciones, Kanban). & SD6, 6.1, Tabla 6.1 y Figura 6.1; T-9, Tabla T-9.1 \\
\LFXfiletefila
\mbox{SD6-03} & Control de cambios sin las reglas del Art. 72°. & \textbf{Corregida.} Plazo de análisis de impacto (5 o 10 días hábiles), tope acumulado con alerta al 15 \% y límite del 20 \%, y regla para cambios en una etapa con fecha fija. & SD6, 6.1.4; T-9, 1.4, Tabla T-9.3 \\
\LFXfiletefila
\mbox{SD6-04} & Escalamiento sin destinatario, plazo ni disparador. & \textbf{Corregida.} Tres niveles con disparadores observables y plazos (mismo día, 2 y 5 días hábiles), y escalamiento inmediato ante amenazas a seguridad, continuidad o hitos; coincide con el SD8. & SD6, 6.1.6, Figura 6.2; T-9, 1.6, Tabla T-9.5 \\
\LFXfiletefila
\mbox{SD6-05} & Matriz de comunicaciones genérica (sin terreno, turno de noche, terceros ni planificador). & \textbf{Corregida.} Matriz con canales de terreno: inicio de jornada en andén, cambio de turno a las 22:00, conductor externo al asignarse, aviso en el terminal, planificador y sindicato. & SD6, 6.1.2; T-9, 1.2, Tabla T-9.2 \\
\LFXfiletefila
\mbox{SD6-06} & Registro de 17 interesados sin influencia ni interés, frente a 19 en SD2 y SD3. & \textbf{Corregida.} Registro de los 19 actores con la influencia e interés del SD2, estrategia por celda, momento, responsable y canal, alineado con el SD3. & SD6-Anexos, Tabla 6.A.1 \\
\LFXfiletefila
\mbox{SD6-07} & Párrafos de evaluadores externos repetidos en SD6 y T-9. & \textbf{Corregida.} Se retiran del SD6 y del T-9; el dimensionamiento queda en el plan de calidad. & SD6, 6.1.3; T-9, 1.3 \\
\LFXfiletefila
\mbox{SD6-08} & Sin método para nivelar la capacidad en los solapamientos 13–15 y 19–20. & \textbf{Corregida.} Control semanal de horas asignadas frente a la curva de plazas del T-15 y el paquete 1.3.4, con reasignación o reserva protegida y escalamiento. & SD6, 6.1.5; T-9, 1.5 \\
\LFXfiletefila
\mbox{SD6-09} & RUP justificado en una frase, sin comparar alternativas. & \textbf{Corregida.} Comparación con cascada, Scrum y SAFe para las condiciones del caso. & SD6, 6.2, Tabla 6.2; T-10, introducción \\
\LFXfiletefila
\mbox{SD6-10} & Requerimientos, arquitectura evolutiva, refactorización y deuda técnica en tres líneas. & \textbf{Corregida.} Gestión de requerimientos (casos de uso, refinamiento con el referente del CLIENTE, trazabilidad y cambios) y tratamiento de la deuda por clase con plazos. & SD6, 6.2.1; T-10, secciones 1 y 2, Tabla T-10.1 \\
\LFXfiletefila
\mbox{SD6-11} & No dice cómo se desarrolla y prueba lo propio de Puelche (desconexión, −22 °C, ERP sin documentar, cadenas, planificador). & \textbf{Corregida.} Seis condiciones del caso con su forma de desarrollo y su prueba antes de la marcha blanca, más la transferencia al área de TI de cuatro personas y el tiempo de salida al mercado. & SD6, 6.2 y 6.2.4, Tabla 6.3; T-10, sección 5 \\
\LFXfiletefila
\mbox{SD6-12} & El T-10 habla de «los cinco ambientes» sin nombrarlos ni definir promoción y aprobación. & \textbf{Corregida.} Tabla de los cinco ambientes con lo que se despliega, la puerta de entrada y quién aprueba. & T-10, Tabla T-10.2; SD6, Figura 6.4 \\
\LFXfiletefila
\mbox{SD6-13} & Iteraciones, ceremonias y releases sin fecha en la planificación. & \textbf{Parcial.} Ceremonias con momento, participantes y artefacto, ligadas al calendario de iteraciones del T-15; la Gantt del SD7 sigue a nivel de paquete. Lo que falta se completa para la Propuesta Final. & SD6, 6.2.3; T-10, Tabla T-10.3; T-15 \\
\LFXfiletefila
\mbox{SD6-14} & Cuerpo de ocho páginas con una sola figura, cuya fuente es equivocada; sin diagramas de gobierno, iteración ni pipeline. & \textbf{Corregida.} Cuerpo de 18 páginas con cuatro figuras (modelo en tres niveles, escalamiento, control de cambios y cadena de entrega) y fuentes corregidas. & SD6, Figuras 6.1 a 6.4 \\
\LFXfiletefila
\mbox{SD6-15} & T-9 y T-10 repiten el SD6 en lugar de detallar. & \textbf{Corregida.} Los formularios agregan tablas de adaptación, procedimiento de cambios, umbrales, escalamiento, contenido de actas, deuda técnica, ambientes y ceremonias. & T-9; T-10 \\
\LFXfiletefila
\mbox{SD6-16} & El SD7 remite a las iteraciones de la 6.2, pero el cuerpo del SD6 no fija su duración (sólo el T-10). & \textbf{Corregida.} El cuerpo fija iteraciones de dos semanas (catorce días) con fecha en el T-15. & SD6, 6.2 y 6.2.3 \\
\LFXfiletefila
\mbox{SD6-17} & Declaración de IA «Alto» en todas las secciones; la revisión declarada es de coherencia, no de autoría. & \textbf{Corregida.} Revisores distribuidos por sección, cada uno con el contenido que verificó. & SD6, Declaración de IA \\
\end{longtable}
}

\Needspace{8\baselineskip}
\subsection{Subdocumento 7: Plan de trabajo, EDT y cronograma (T-14, T-15 y T-18)}
{\linespread{1}\LFXcuerpotabla
\setlength{\parskip}{0pt}\renewcommand{\arraystretch}{1.10}\setlength{\tabcolsep}{4pt}
\begin{longtable}{F{0.09} F{0.29} F{0.40} F{0.22}}
\caption[Subdocumento 7: Plan de trabajo, EDT y cronograma (T-14, T-15 y T-18)]{Subdocumento 7: Plan de trabajo, EDT y cronograma (T-14, T-15 y T-18). Fuente: revisión del Informe 2 y versión corregida de LafroX.}\\[2pt]
\LFXtopruleoriginal \cab{N.º} & \cab{Observación} & \cab{Respuesta} & \cab{Sección modificada} \\ \midrule
\endfirsthead
\multicolumn{4}{@{}l@{}}{\fontsize{9pt}{11pt}\selectfont\color{lafroxGris}Tabla \thetable\ (continuación)}\\[3pt]
\LFXtopruleoriginal \cab{N.º} & \cab{Observación} & \cab{Respuesta} & \cab{Sección modificada} \\ \midrule
\endhead
\midrule
\multicolumn{4}{@{}r@{}}{\fontsize{9pt}{11pt}\selectfont\color{lafroxGris}Continúa en el folio siguiente.}\\
\endfoot
\bottomrule
\endlastfoot
\mbox{SD7-01} & Primer nivel de la EDT organizado por fases RUP y no por entregables. & \textbf{Parcial.} Se agrega en el T-14 un índice de la EDT por productos verificables; las nueve ramas por fase se conservan. Lo que falta se completa para la Propuesta Final. & SD7, 7.1; T-14, índice por productos \\
\LFXfiletefila
\mbox{SD7-02} & Etapas literales dentro de las ramas (3.4/3.5, 3.8/3.9, 4.2/4.3). & \textbf{Parcial.} La etapa se declara como incremento que habilita cada producto; las ramas se mantienen. Lo que falta se completa para la Propuesta Final. & SD7, 7.1 \\
\LFXfiletefila
\mbox{SD7-03} & Estimación por clases fijas (todo módulo vale 960 HH). & \textbf{Corregida.} Estimación por tamaño de cada módulo: flujos S/M/C, interfaces y dispositivos trazables al T-12, con actividades de diseño, construcción, pruebas e integración. & T-15, 4.2 (estimación por tamaño) \\
\LFXfiletefila
\mbox{SD7-04} & PERT degenerado (O = 0,75M y P = 1,25M para todos). & \textbf{Corregida.} O y P con fundamento propio por módulo, tipo de paquete y riesgo; la banda del 25 \% queda sólo como sensibilidad. & T-15, 4.2 y 5 \\
\LFXfiletefila
\mbox{SD7-05} & Capacidad de calidad insuficiente (faltan 217,60 HH) e implantación sin división que la provea. & \textbf{Corregida.} Tres refuerzos de pruebas (1.081,60 HH frente a 1.049,60), área de Implantación con 9 personas y 18 asistentes por obra; 136 plazas conciliadas con el SD1. & T-15, 5.7 y 6.3; SD1, 1.2 \\
\LFXfiletefila
\mbox{SD7-06} & Once módulos de la Etapa 1 en paralelo, sobre el techo de desarrolladores del SD1. & \textbf{Corregida.} La curva de plazas por grupo de origen queda dentro de la planta: Desarrollo llega a 32 plazas en los meses 6 y 7. & T-15, 4.4 y 5.7 \\
\LFXfiletefila
\mbox{SD7-07} & «Incompatibilidad por resolver» entre la revisión del acta (Art. 18.3) y el paso a producción. & \textbf{Corregida.} Acta integrada por etapa con diez días hábiles de revisión; producción el 17-05-2028 (H7) y el 18-10-2028 (H12), con 10 y 8 días hábiles de holgura. & T-15, 5.5 y Tabla 7; SD7, 7.3 \\
\LFXfiletefila
\mbox{SD7-08} & Los once módulos de la Etapa 1 arrancan la misma semana; la secuencia por dependencias que describe el texto no se modela como orden de construcción. & \textbf{Pendiente.} Los paquetes 3.4.1 a 3.4.10 siguen comenzando entre el 10 y el 18 de junio de 2027. Se corrige para la Propuesta Final. & T-15, 6.1 (programación de paquetes) \\
\LFXfiletefila
\mbox{SD7-09} & Red incompleta: 102 de 168 paquetes sin predecesora (migración, cadenas, innovaciones). & \textbf{Parcial.} Se agregan predecesoras; quedan 77 paquetes sin vínculo declarado. Lo que falta se completa para la Propuesta Final. & T-15, 6.1 (programación de paquetes) \\
\LFXfiletefila
\mbox{SD7-10} & CPM calculado sobre fechas publicadas y no sobre duraciones. & \textbf{Corregida.} El CPM se calcula desde las duraciones y las predecesoras de los 227 paquetes. & T-15, 5 \\
\LFXfiletefila
\mbox{SD7-11} & Iteraciones de dos semanas sólo como plantilla; no se ven en la Gantt de construcción. & \textbf{Parcial.} Calendario de iteraciones por módulo con demostración y regresión; la Gantt sigue a nivel de paquete. Lo que falta se completa para la Propuesta Final. & T-15, calendario de iteraciones \\
\LFXfiletefila
\mbox{SD7-12} & Pruebas concentradas después de construir (cierre en V). & \textbf{Parcial.} Cada iteración cierra con pruebas del incremento en QA; las pruebas ampliadas 3.8 y 3.9 mantienen su ventana. Lo que falta se completa para la Propuesta Final. & T-15, 6; SD9, 9.3 \\
\LFXfiletefila
\mbox{SD7-13} & Gantt con rótulos de 5 a 6 pt, encabezados superpuestos y códigos sin definir. & \textbf{Corregida.} Láminas vectoriales con nombres completos, inicio y fin, eje temporal y leyenda (mínimo 9,5 pt). & SD7, Figuras 7.24 a 7.31 \\
\LFXfiletefila
\mbox{SD7-14} & Filas agrupadas que mezclan paquetes de cuentas distintas. & \textbf{Parcial.} Códigos completos y hojas de relaciones por ID; las agrupaciones del modelo se conservan. Lo que falta se completa para la Propuesta Final. & SD7, Figura 7.26 \\
\LFXfiletefila
\mbox{SD7-15} & Defecto crítico durante el congelamiento de septiembre en la marcha blanca de la Etapa 2. & \textbf{Corregida.} Procedimiento: registro de severidad, aislamiento con medidas autorizadas y escalamiento a la Contraparte Técnica, sin cambios en Producción. & SD7, 7.3; T-18 \\
\LFXfiletefila
\mbox{SD7-16} & Reversión: 10 min en la Tabla 7.8 y 40 min en el 7.3.8 sin explicación. & \textbf{Corregida.} Se explican los niveles: max(10, 30) + 10 = 40 min, con decisión a las 04:45 y restauración a las 05:25. & SD7, 7.3, Tabla 7.8 \\
\LFXfiletefila
\mbox{SD7-17} & El intercambio con la principal cadena parte en el mes 13, sin margen ante agosto de 2028. & \textbf{Pendiente.} La ventana del paquete 3.6.5 se adelanta en la Propuesta Final para dejar margen antes de agosto de 2028. & T-14, 3.6.5 \\
\LFXfiletefila
\mbox{SD7-18} & Sin evidencia al proveedor de lácteos entre el inicio y el mes 12. & \textbf{Pendiente.} Se agrega en la Propuesta Final una entrega de evidencia al proveedor de lácteos antes del mes 12. & T-14, 1.1.3 \\
\LFXfiletefila
\mbox{SD7-19} & Errata del 1.2.1: «Mes 1 (H1)» con la línea base firmada en el mes 2. & \textbf{Corregida.} Elaboración en el mes 1 y aprobación en H1 (mes 2); igual corrección en 1.2.5 (H8). & T-14, Tabla 4 \\
\LFXfiletefila
\mbox{SD7-20} & Tablas 7.2 y 7.8 con párrafos en las celdas. & \textbf{Corregida.} Celdas breves y explicación en prosa contigua. & SD7, Tablas 7.2 y 7.8 \\
\LFXfiletefila
\mbox{SD7-21} & El 7.3.8 arrastra contenido ajeno al plan, duplicado en el T-18. & \textbf{Parcial.} Se condensa el 7.3.8 con remisión al T-18, cuyo bloque cambia de título. Lo que falta se completa para la Propuesta Final. & SD7, 7.3.8; T-18, 6.5 \\
\LFXfiletefila
\mbox{SD7-22} & Siglas SD4, SD8, PT/H y RC sin definir; Figura 7.20 remite a un detalle inexistente. & \textbf{Corregida.} Siglas definidas en los cinco documentos y leyendas; se elimina la llamada de la Figura 7.20. & SD7, T-14, T-15, T-18; Figura 7.20 \\
\LFXfiletefila
\mbox{SD7-23} & Volumen: 66 páginas de cuerpo y 250 de formularios y anexos. & \textbf{Pendiente.} La reducción del cuerpo y de los formularios se aplica en la Propuesta Final; esta versión priorizó la legibilidad de las figuras. & SD7 y formularios \\
\LFXfiletefila
\mbox{SD7-24} & Redacción defensiva («no acredita», «no demuestra», «no se presume»). & \textbf{Parcial.} Depuración parcial en el T-18 y el 7.3.8. Lo que falta se completa para la Propuesta Final. & SD7; T-15; T-18 \\
\LFXfiletefila
\mbox{SD7-25} & Falta la ruta crítica de la construcción de la Etapa 1, no sólo la del tramo final. & \textbf{Parcial.} La tabla de CPM calcula todas las ramas, incluida la construcción; la figura sigue mostrando sólo la cadena que controla H12. Lo que falta se completa para la Propuesta Final. & T-15, 5 (ruta crítica y holguras) \\
\LFXfiletefila
\mbox{SD7-26} & Sólo las 119 actividades de los paquetes que terminan antes del H2 están detalladas; el resto queda a planificación gradual. & \textbf{Pendiente.} Se mantiene la planificación gradual después del H2. Se corrige para la Propuesta Final. & T-15, 6 \\
\end{longtable}
}

\Needspace{8\baselineskip}
\subsection{Subdocumento 8: Plan de riesgos (T-16)}
{\linespread{1}\LFXcuerpotabla
\setlength{\parskip}{0pt}\renewcommand{\arraystretch}{1.10}\setlength{\tabcolsep}{4pt}
\begin{longtable}{F{0.09} F{0.29} F{0.40} F{0.22}}
\caption[Subdocumento 8: Plan de riesgos (T-16)]{Subdocumento 8: Plan de riesgos (T-16). Fuente: revisión del Informe 2 y versión corregida de LafroX.}\\[2pt]
\LFXtopruleoriginal \cab{N.º} & \cab{Observación} & \cab{Respuesta} & \cab{Sección modificada} \\ \midrule
\endfirsthead
\multicolumn{4}{@{}l@{}}{\fontsize{9pt}{11pt}\selectfont\color{lafroxGris}Tabla \thetable\ (continuación)}\\[3pt]
\LFXtopruleoriginal \cab{N.º} & \cab{Observación} & \cab{Respuesta} & \cab{Sección modificada} \\ \midrule
\endhead
\midrule
\multicolumn{4}{@{}r@{}}{\fontsize{9pt}{11pt}\selectfont\color{lafroxGris}Continúa en el folio siguiente.}\\
\endfoot
\bottomrule
\endlastfoot
\mbox{SD8-01} & 13 riesgos internos del proveedor declarados al CLIENTE. & \textbf{Parcial.} Se retiran del registro como controles internos (CI) con solución y cierre; el T-16 queda con 35 riesgos. Falta reflejarlos en SD4, SD9 y SD13. Lo que falta se completa para la Propuesta Final. & SD8-Anexos, 8.E.1 y 8.I.2; T-16 \\
\LFXfiletefila
\mbox{SD8-02} & Riesgos del CLIENTE sin contraparte, pedido ni fecha. & \textbf{Corregida.} Cada riesgo nombra la contraparte, qué se pide, el plazo y la contingencia si no responde. & SD8-Anexos, 8.A.2 y 8.F.1; T-16 \\
\LFXfiletefila
\mbox{SD8-03} & No hay riesgo residual. & \textbf{Corregida.} Valoración posterior de los 35 riesgos con probabilidad, impacto, zona e hipótesis. & SD8-Anexos, 8.F.2; T-16 \\
\LFXfiletefila
\mbox{SD8-04} & Riesgos en zona Alta o Crítica sin plan adicional. & \textbf{Parcial.} Plan adicional con plazo, dueño y bloqueo; no se demuestra aún un nivel manejable en todos. Lo que falta se completa para la Propuesta Final. & SD8-Anexos, 8.F.1, 8.F.2 y 8.G.1 \\
\LFXfiletefila
\mbox{SD8-05} & Faltan riesgos del Cap. 19 (sindicato, rotación, interfaces ERP/WMS, Concepción, peak) y obsolescencia, bloqueo y escalabilidad. & \textbf{Parcial.} Se incorporan como R-028 a R-035 (incluido el rechazo de conductores externos); la escalabilidad se trata como control interno. Lo que falta se completa para la Propuesta Final. & T-16; SD8-Anexos, 8.D.3 y 8.E.1 \\
\LFXfiletefila
\mbox{SD8-06} & Puntos únicos de falla del SD4 no llevados al registro. & \textbf{Parcial.} El ERP queda en el T-16; identidad y stock como controles internos. Falta sincronizar con SD4 y SD9. Lo que falta se completa para la Propuesta Final. & SD8-Anexos, 8.D y 8.E.1 \\
\LFXfiletefila
\mbox{SD8-07} & No se aplica FMEA, árbol de fallas ni simulación. & \textbf{Corregida.} FMEA de 15 modos con S/O/D, NPR, controles y puntuación posterior. & SD8, 8.2.3; Anexo 8.G.1 \\
\LFXfiletefila
\mbox{SD8-08} & Sin análisis costo-beneficio. & \textbf{Corregida.} Quince cálculos en HH con sensibilidad y justificación de los retornos menores que 1. & SD8-Anexos, 8.F.1 y 8.F.3 \\
\LFXfiletefila
\mbox{SD8-09} & Acciones sin plazo, esfuerzo ni contingencia. & \textbf{Corregida.} Treinta y cinco fichas con perfiles, HH y actuación contingente temporizada. & SD8-Anexos, 8.F.1 \\
\LFXfiletefila
\mbox{SD8-10} & Reservas sin derivar ni reflejo en el cronograma; 3.072 HH del SD7 no conciliadas. & \textbf{Parcial.} Reserva de gestión derivada por escenario y conciliada con la reserva protegida; falta incorporarla al T-15. Lo que falta se completa para la Propuesta Final. & SD8, 8.3.2; Anexo 8.H \\
\LFXfiletefila
\mbox{SD8-11} & Escala de impacto genérica, factores sin definir, bandas solapadas en 70 \%, frase que niega los umbrales de zona y dos probabilidades por riesgo. & \textbf{Corregida.} Factores, bandas y umbrales únicos definidos en el cuerpo; se retira la frase sobre la zona de prioridad. & SD8, 8.1.3; Tablas 8.2 a 8.4 \\
\LFXfiletefila
\mbox{SD8-12} & Dos «Tabla 8.4», Tabla 8.5 de siete columnas, fuente errada del RT-20 y RBS sin folio. & \textbf{Corregida.} Numeración única, tablas de cinco columnas, fuente corregida y RBS nueva con folio y firma. & SD8, 8.1 a 8.3; Figura 8.1 \\
\LFXfiletefila
\mbox{SD8-13} & Fuentes de las Tablas 8.2 a 8.4 y de la Declaración de IA con «contenido academico». & \textbf{Corregida.} Las fuentes pasan a «elaboración propia» y la Declaración de IA se reescribe sin esa mención. & SD8, Tablas 8.2 a 8.4 y Declaración de IA \\
\LFXfiletefila
\mbox{SD8-14} & Erratas: «tendria», «basandose», «academico», «asegurandose», «estmación». & \textbf{Corregida.} Erratas corregidas en el cuerpo, el T-16 y los anexos. & SD8; T-16; SD8-Anexos \\
\LFXfiletefila
\mbox{SD8-15} & El riesgo del planificador no menciona la jubilación con fecha conocida ni la caída de 6 a 9 puntos. & \textbf{Corregida.} R-014 trata la jubilación dentro del horizonte contractual y la pérdida de 6 a 9 puntos de cumplimiento. & T-16, R-014 \\
\LFXfiletefila
\mbox{SD8-16} & El SD13 asigna al piloto de R-034 meses que el Anexo 8.A no indica. & \textbf{Parcial.} CI-034 fija paquete 11–15 y piloto 13–15; falta actualizar la referencia en el SD13. Lo que falta se completa para la Propuesta Final. & SD8-Anexos, 8.E.1 \\
\LFXfiletefila
\mbox{SD8-17} & Falla mecánica, siniestro y asalto a camión: eventos operacionales del cliente de poco valor para el plan. & \textbf{Pendiente.} Se mantienen como R-021 a R-023 por su efecto en la continuidad del reparto; su tratamiento se acota en la Propuesta Final. & T-16 \\
\LFXfiletefila
\mbox{SD8-18} & El T-16 menciona un «VME» que nunca se define ni se calcula. & \textbf{Corregida.} Se retira la mención. & T-16 \\
\LFXfiletefila
\mbox{SD8-19} & La Declaración de IA dice que la IA generó los riesgos y la tabla del T-16, con un solo revisor. & \textbf{Corregida.} La fila del T-16 se reformula (la herramienta mejora la redacción y propone valores) y cinco revisores se reparten las secciones, cada uno con el contenido que verificó. & SD8, Declaración de IA \\
\LFXfiletefila
\mbox{SD8-20} & Estilo defensivo («no certifican», «no se presume», «no se afirma»). & \textbf{Parcial.} Se reducen esas expresiones de 13 a 10 en el cuerpo y el T-16. Lo que falta se completa para la Propuesta Final. & SD8; T-16 \\
\LFXfiletefila
\mbox{SD8-21} & Cuerpo de unas 14 páginas de contenido para un ítem del 10 \%. & \textbf{Pendiente.} El cuerpo mantiene su extensión; el desarrollo nuevo (controles, respuestas, FMEA y reservas) va en los Anexos 8.E a 8.H. Se corrige para la Propuesta Final. & SD8; SD8-Anexos, 8.E a 8.H \\
\end{longtable}
}

\Needspace{8\baselineskip}
\subsection{Subdocumento 9: Plan de calidad (T-13 y T-17)}
{\linespread{1}\LFXcuerpotabla
\setlength{\parskip}{0pt}\renewcommand{\arraystretch}{1.10}\setlength{\tabcolsep}{4pt}
\begin{longtable}{F{0.09} F{0.29} F{0.40} F{0.22}}
\caption[Subdocumento 9: Plan de calidad (T-13 y T-17)]{Subdocumento 9: Plan de calidad (T-13 y T-17). Fuente: revisión del Informe 2 y versión corregida de LafroX.}\\[2pt]
\LFXtopruleoriginal \cab{N.º} & \cab{Observación} & \cab{Respuesta} & \cab{Sección modificada} \\ \midrule
\endfirsthead
\multicolumn{4}{@{}l@{}}{\fontsize{9pt}{11pt}\selectfont\color{lafroxGris}Tabla \thetable\ (continuación)}\\[3pt]
\LFXtopruleoriginal \cab{N.º} & \cab{Observación} & \cab{Respuesta} & \cab{Sección modificada} \\ \midrule
\endhead
\midrule
\multicolumn{4}{@{}r@{}}{\fontsize{9pt}{11pt}\selectfont\color{lafroxGris}Continúa en el folio siguiente.}\\
\endfoot
\bottomrule
\endlastfoot
\mbox{SD9-01} & Falta la prueba del manifiesto de identidad durante un corte. & \textbf{Corregida.} Protocolo AL-ID-01 (emisión, relevo, expiración y revocación sin enlace) con criterio de salida. & SD9, 9.2, Tabla 9.8; T-13 \\
\LFXfiletefila
\mbox{SD9-02} & Usabilidad sin población ni horario (rotación del 38 \%). & \textbf{Corregida.} Prueba en turno nocturno con personal de menos de un mes; aprendizaje ≤ 2 h, error ≤ 5 \% y criterio semanal. & SD9, 9.1 y 9.2; T-13 \\
\LFXfiletefila
\mbox{SD9-03} & Dos términos para el paso a producción y dos reglas para cambios en congelamiento. & \textbf{Corregida.} Se distingue marcha blanca de producción como en el T-18 y queda una sola regla de congelamiento. & SD9, 9.2 y 9.3; T-13 \\
\LFXfiletefila
\mbox{SD9-04} & Evaluación CMMI con renovación previa a la adjudicación. & \textbf{Corregida.} Se presenta como plan de renovación pendiente, sin invocarla como garantía. & SD9, 9.1 \\
\LFXfiletefila
\mbox{SD9-05} & Restauración mensual en septiembre y diciembre sin ambiente declarado. & \textbf{Corregida.} Se especifica la restauración de prueba durante el congelamiento. & SD9, 9.3 \\
\LFXfiletefila
\mbox{SD9-06} & Errata «extbfCorrecciones en congelamiento» y rótulo superpuesto en la Figura 9.4. & \textbf{Corregida.} Encabezado corregido y rótulo separado del recuadro. & SD9, 9.3; Figura 9.4 \\
\LFXfiletefila
\mbox{SD9-07} & Salvedades defensivas que oscurecen las reglas. & \textbf{Corregida.} Se simplifican sin cambiar las reglas. & SD9, 9.1 a 9.3 \\
\LFXfiletefila
\mbox{SD9-08} & Pruebas concentradas al final y holgura negativa en H1, H7 y H12. & \textbf{Parcial.} Pruebas por quincena y holguras contrastadas con el T-15; depende de la Gantt del SD7. Lo que falta se completa para la Propuesta Final. & SD9, 9.3, Tabla 9.13 \\
\LFXfiletefila
\mbox{SD9-09} & Matriz del T-13 de 50 páginas sin resumen. & \textbf{Corregida.} Resumen por módulo al inicio de la sección. & T-13, 8.2 (resumen por módulo) \\
\end{longtable}
}

\Needspace{8\baselineskip}
\subsection{Subdocumento 13: Innovaciones (T-19)}
{\linespread{1}\LFXcuerpotabla
\setlength{\parskip}{0pt}\renewcommand{\arraystretch}{1.10}\setlength{\tabcolsep}{4pt}
\begin{longtable}{F{0.09} F{0.29} F{0.40} F{0.22}}
\caption[Subdocumento 13: Innovaciones (T-19)]{Subdocumento 13: Innovaciones (T-19). Fuente: revisión del Informe 2 y versión corregida de LafroX.}\\[2pt]
\LFXtopruleoriginal \cab{N.º} & \cab{Observación} & \cab{Respuesta} & \cab{Sección modificada} \\ \midrule
\endfirsthead
\multicolumn{4}{@{}l@{}}{\fontsize{9pt}{11pt}\selectfont\color{lafroxGris}Tabla \thetable\ (continuación)}\\[3pt]
\LFXtopruleoriginal \cab{N.º} & \cab{Observación} & \cab{Respuesta} & \cab{Sección modificada} \\ \midrule
\endhead
\midrule
\multicolumn{4}{@{}r@{}}{\fontsize{9pt}{11pt}\selectfont\color{lafroxGris}Continúa en el folio siguiente.}\\
\endfoot
\bottomrule
\endlastfoot
\mbox{SD13-01} & Ninguna innovación cuantifica el beneficio en unidades. & \textbf{Parcial.} Cada innovación dimensiona su beneficio en unidades físicas, con escenarios ilustrativos sobre supuestos declarados. Lo que falta se completa para la Propuesta Final. & SD13, 13.1.5 a 13.5.5; T-19 \\
\LFXfiletefila
\mbox{SD13-02} & INN-03 cuenta como innovación y como cumplimiento del RT-05.30. & \textbf{Corregida.} Se separa el modelo analítico ofertado (RT-05.30) del aporte propio de la innovación. & SD13, 13.3 \\
\LFXfiletefila
\mbox{SD13-03} & INN-02 es una práctica del proveedor más que un proceso del CLIENTE. & \textbf{Parcial.} El beneficio para Puelche se expresa en minutos de interrupción evitados en terreno. Lo que falta se completa para la Propuesta Final. & SD13, 13.2 \\
\LFXfiletefila
\mbox{SD13-04} & La valorización «se desarrolla para el Informe 3». & \textbf{Corregida.} Se remite a la Oferta Económica, Entregable 2. & SD13, introducción y 13.x.5 \\
\LFXfiletefila
\mbox{SD13-05} & Revisión humana sólo de redacción; sin revisión técnica. & \textbf{Parcial.} Pauta de revisión técnica de modelos, umbrales y cálculos por innovación; su aplicación completa se registra en la Propuesta Final. & SD13-Anexos, pauta de revisión técnica \\
\LFXfiletefila
\mbox{SD13-06} & Fechas de paquetes por verificar con los capítulos nuevos. & \textbf{Corregida.} Validaciones y holguras conciliadas con el T-15; controles y riesgos alineados con el SD8. & SD13; T-19 \\
\end{longtable}
}


\microtypesetup{protrusion=false}
\Needspace{6\baselineskip}
\section*{Referencias}
\addcontentsline{toc}{section}{Referencias}
\begin{itemize}[itemsep=1pt]
\item Distribuidora Puelche S.A. (2026a). \emph{Bases Administrativas de Licitación N.º TFEP-01/2026: Contratación de Solución Integral de Software y Servicios de Operación}.
\item Distribuidora Puelche S.A. (2026b). \emph{Bases Técnicas Transversales de Licitación N.º TFEP-01/2026}.
\item Distribuidora Puelche S.A. (2026c). \emph{Bases Técnicas del caso: Caso 02 --- Logística}.
\end{itemize}

\clearpage
\section*{Declaración de uso de IA}
\addcontentsline{toc}{section}{Declaración de uso de IA}

La Tabla~\ref{tab:declaracion_ia} identifica la herramienta y el alcance de la asistencia utilizada en este anexo.

\begin{tablalafrox}{Uso de IA en el Anexo A.}{tab:declaracion_ia}{L{1.6cm} L{2.0cm} Y L{1.2cm} L{1.6cm} L{3.4cm}}
  {\cab{Sección} & \cab{Herramienta} & \cab{Finalidad del uso} & \cab{Nivel en texto} & \cab{Nivel en diagramas} & \cab{Revisión humana (quién y qué verificó)}}
  A.1 & Claude Code & Redacción del alcance y de la síntesis & Alto & Ninguno & Leandro Chamorro: criterio de registro, conteo de estados y síntesis por subdocumento. \\
  A.2 & Claude Code & Cotejo de cada observación con la versión corregida de los subdocumentos y redacción de las respuestas & Alto & Ninguno & Álvaro Catalán (Subdocumentos 1 a 5) y Tomás Pérez (Subdocumentos 6 a 9 y 13): estado y sección modificada de cada fila contra el subdocumento correspondiente. \\
\end{tablalafrox}

\end{document}
